\chapter{Recursion}

\section{What Recursion Means}\index{recursion}
A \index{recursion}\newterm{recursive} function is a function that calls itself. Think of Russian nesting dolls. To find the smallest doll, you open the big doll, then open the next one, and so on, until you reach a doll that does not open.
Many problems can be split into a smaller copy of the same problem. Recursion
solves the big problem by solving the smaller one and then combining the answers.

Every recursive function needs two parts:
\begin{itemize}
\item A \index{base case}\newterm{base case}, which returns an answer directly
without calling the function again.
\item A \newterm{recursive case}, which calls the function on a smaller input.
\end{itemize}
If the base case is missing, or if the recursive case never gets closer to it,
the function calls itself forever.

\section{Factorial}
The factorial $n!$ is $n \cdot (n - 1)!$, with $0! = 1$. That definition is
recursive, and the code is almost the same:

\begin{minted}{python}
def factorial(n):
    if n == 0:
        return 1
    return n * factorial(n - 1)

print(factorial(5))  # 120
\end{minted}

The base case is \pyvar{n == 0}. Each call makes the input one smaller, so it
reaches the base case after $n$ calls.

\section{The Call Stack}\index{call stack}
Each time a function is called, Python stores a \newterm{frame} on the call
stack. The frame holds that call's variables. When the call returns, its frame
is removed. For \pyfunction{factorial(3)}, the stack grows to four frames:
\texttt{factorial(3)}, \texttt{factorial(2)}, \texttt{factorial(1)}, and
\texttt{factorial(0)}. Then the answers return back up the stack: $1$, $1$, $2$, and $6$.

\section{Binary Search, Recursively}\index{binary search!recursive}
\newterm{Binary search} finds an item in a \emph{sorted} list the way you
look up a word in a dictionary: open to the middle, decide which half the
word must be in, and repeat on that half. Each step throws away half of
what is left, so even a million items take only about 20 steps. It is
naturally recursive. Look at the middle item. If it matches, return its
index. Otherwise, search only the half that could contain the target.

\begin{minted}{python}
def binary_search(items, target, low, high):
    if low > high:
        return -1
    mid = (low + high) // 2
    if items[mid] == target:
        return mid
    elif items[mid] < target:
        return binary_search(items, target, mid + 1, high)
    else:
        return binary_search(items, target, low, mid - 1)

data = [1, 3, 5, 7, 9, 11]
print(binary_search(data, 7, 0, len(data) - 1))  # 3
print(binary_search(data, 4, 0, len(data) - 1))  # -1
\end{minted}

The first base case, \pyvar{low > high}, handles a target that is not in the
list. The second, \pyvar{items[mid] == target}, handles a match.

\section{Recursive Traversal of a Tree}\index{depth-first search!recursive}
A \emph{tree} is a structure like a family tree or the folders on your
computer: each node can have children, which can have children of their
own. (You will study trees in detail in the next chapter.) A tree is
naturally recursive: each child is the root of a smaller tree. To visit every node, visit the node, then visit each of its
subtrees.

\begin{minted}{python}
class Node:
    def __init__(self, name, children=None):
        self.name = name
        self.children = children or []

def visit(node):
    names = [node.name]
    for child in node.children:
        names.extend(visit(child))
    return names

tree = Node('a', [Node('b', [Node('d')]), Node('c')])
print(visit(tree))  # ['a', 'b', 'd', 'c']
\end{minted}

A node with no children is the base case. Its loop does nothing, so it
returns just its own name.

\section{The Cost of Repeated Work}\index{Fibonacci numbers}
The Fibonacci numbers are defined by $F(n) = F(n-1) + F(n-2)$, with $F(0) = 0$
and $F(1) = 1$. A direct recursive version is short:

\begin{minted}{python}
def fib(n):
    if n < 2:
        return n
    return fib(n - 1) + fib(n - 2)
\end{minted}

It is correct, but it is slow. The call for \texttt{fib(10)} makes 177 calls in
total, and \texttt{fib(15)} makes 1,973. The same smaller values are computed
again and again. The number of calls grows exponentially with $n$.

\newterm{Memoization} fixes this by saving each answer the first time it is
computed:

\begin{minted}{python}
def fib_memo(n, memo=None):
    if memo is None:
        memo = {}
    if n < 2:
        return n
    if n not in memo:
        memo[n] = fib_memo(n - 1, memo) + fib_memo(n - 2, memo)
    return memo[n]

print(fib_memo(80))  # 23416728348467685
\end{minted}

\begin{Exercise}[label=rec1, title=Trace the calls]
Write out every call made by \texttt{factorial(3)}, in order, and then the value
returned by each call.
\end{Exercise}
\begin{Answer}[ref=rec1]
Calls: \texttt{factorial(3)}, \texttt{factorial(2)}, \texttt{factorial(1)},
\texttt{factorial(0)}. Returns: $1$, then $1 \cdot 1 = 1$, then $2 \cdot 1 = 2$,
then $3 \cdot 2 = 6$.
\end{Answer}

\begin{Exercise}[label=rec2, title=Sum a list]
Write a recursive function \pyfunction{total(items)} that returns the sum of a list.
What is its base case?
\end{Exercise}
\begin{Answer}[ref=rec2]
\begin{minted}{python}
def total(items):
    if len(items) == 0:
        return 0
    return items[0] + total(items[1:])
\end{minted}
The base case is the empty list, whose sum is 0.
\end{Answer}

\begin{Exercise}[label=rec4, title=Counting calls]
How many calls does \texttt{fib(5)} make in total? Use the formula that the number of calls for \texttt{fib(n)} is $2F(n+1) - 1$, where $F$ is the Fibonacci sequence, and check the answer by counting the calls.
\end{Exercise}
\begin{Answer}[ref=rec4]
$F(6) = 8$, so the formula gives $2 \cdot 8 - 1 = 15$ calls. Counting: \texttt{fib(5)} calls \texttt{fib(4)} and \texttt{fib(3)}, and so on. The total is 15, which matches the formula.
\end{Answer}

\begin{Exercise}[label=rec3, title=Missing base case]
What happens when \pyfunction{factorial(-1)} is called, using the code above?
\end{Exercise}
\begin{Answer}[ref=rec3]
The value $-1$ never reaches the base case $n = 0$. Each call goes to $-2$, $-3$,
and so on, until Python raises a \texttt{RecursionError}.
\end{Answer}
